\documentclass[10pt,a4paper]{amsart}
\usepackage[margin=19mm]{geometry}
\usepackage{amsmath,amssymb,mathtools}
\usepackage[scr=rsfs,cal=euler]{mathalfa}
\usepackage{xfrac}
\usepackage[unicode,hidelinks,bookmarks=false]{hyperref}
\newcommand{\ie}{i.e.\ }
\newcommand{\cf}{cf.\ }
\newcommand{\resp}{respectively\ }
\newcommand{\Subsection}{\subsection}
\providecommand{\nobreakdash}{\nobreak\hbox{-}}
\setlength{\emergencystretch}{2em}
\multlinegap0pt
\usepackage{amssymb}
\usepackage{tikz}
\newtheorem{prop}{Proposition}
\newtheorem{claim}{Claim}
\DeclareMathOperator{\ev}{ev}
\newcommand{\ol}{\overline}
\newcommand{\ra}{\rightarrow}
\title{Rational curves on cubic hypersurfaces in positive characteristic}
\author{Natsume Kitagawa}
\date{Source: arXiv:2604.26556v2 (CC BY 4.0)}
\begin{document}
\maketitle

\noindent\emph{Source and licence.} Rational curves on cubic hypersurfaces in positive characteristic. Version arXiv:2604.26556v2 (CC BY 4.0), distributed by arXiv under the Creative Commons Attribution 4.0 International licence, http://creativecommons.org/licenses/by/4.0/, as recorded in the arXiv licence metadata for this exact version and shown on the abstract page https://arxiv.org/abs/2604.26556v2. Full licence notice: \texttt{licenses/arxiv-2604.26556-CC-BY-4.0.txt}.\par\smallskip

\noindent\textit{Editorial scope.} Counted unit: the article's claim claim1 (source0.tex lines 349--351). The article's complete original proof of it is delivered with it (source0.tex lines 352--357, one \texttt{proof} environment of 6 lines). No mathematics is written, repaired or re-derived: the statement and the proof are the article's own lines, in the article's own notation, and no retained line is substituted or re-flowed.  Retained uncounted as the source-local context the proof needs: lines 297--306; lines 312--314.  Nothing was omitted from any retained span, so the package carries no omission record.  No retained line contains a citation, so the wrapper carries no bibliography and no bibliography entry.  No retained line contains a figure environment, an image-inclusion command or drawing code, so no image is delivered and the package declares no asset dependency.  Editorial formatting only, in addition: an AMS \texttt{amsart} preamble that restates the article's own declarations and macros used by the retained lines, namely the environments \texttt{prop}, \texttt{claim} on separate counters, together with the standard packages those lines need.  The retained environments print in this document as Proposition 1, Claim 1 in order of appearance, whereas the article prints the same items with the numbering of its own full document.  The complete native source of the article as distributed in the arXiv e-print for this exact version is bundled unchanged as \texttt{sources/199404/source0.tex}.
	\begin{prop}\label{DimOfFibres}
		Let $X$ be a smooth cubic hypersurface of dimension $n\geq3$. For an integer $d\geq1$, let $\ev_d:\ol{M}_{0,1}(X,d)\ra X$ be the evaluation map. Then there exists a finite set $K\subset X$ such that
		\begin{align*}
			\dim\ev_d^{-1}(x)=
			\begin{cases}
				(n-1)d-2\qquad\text{if}\quad x\notin K,\\
				(n-1)d-1\qquad\text{if}\quad x\in K.
			\end{cases}
		\end{align*}
	\end{prop}

	\begin{prop}\label{IrredDomain}
		Let $X$ be a smooth cubic hypersurface of dimension $n\geq4$ and let $K\subset X$ be a subset as in Proposition \ref{DimOfFibres}. If $p\notin K$, the fibre $\ev_d^{-1}(p)$ generically parametrizes curves with irreducible domains.
	\end{prop}

		\begin{claim}\label{claim1}
			The main component of $\pi^{-1}(\ev_e^{-1}(x))\times_X\ol{M}_{0,1}(X,d-e)$ is unique. Moreover, $\ev_d^{-1}(x)\cap\Delta_{\{e,\{1\}\},\{d-e,\emptyset\}}$ is irreducible.
		\end{claim}

			\begin{proof}
				The general fibre of the first projection
				$$\pi_1:\pi^{-1}(\ev_e^{-1}(x))\times_X\ol{M}_{0,1}(X,d-e)\ra\pi^{-1}(\ev_e^{-1}(x))$$
				is irreducible by the induction hypothesis. This means the dominant component of $\pi_1$ is unique.
				By Proposition \ref{IrredDomain}, the general map in $\ev_e^{-1}(x)$ has irreducible domain. Hence the component of $\pi^{-1}(\ev_e^{-1}(x))$ which has the maximal dimension is unique, since $\ev_e^{-1}(x)$ is irreducible by the induction hypothesis. Consider the restriction of the glueing map to $\pi^{-1}(\ev_e^{-1}(x))\times_X\ol{M}_{0,1}(X,d-e)$. Note that its image coincides with $\ev_d^{-1}(x)\cap\Delta_{\{e,\{1\}\},\{d-e,\emptyset\}}$. By Proposition \ref{DimOfFibres}, we conclude that $\ev_d^{-1}(x)\cap\Delta_{\{e,\{1\}\},\{d-e,\emptyset\}}$ is irreducible.
			\end{proof}

\end{document}
